%!TEX root = main.tex

The ExpandConnects transformation expands the connect statements between aggregate
types into those between ground types. Correspondingly, it also flattens the aggregate-type components to a collection of ground-type components, so that after the ExpandConnects transformation, all the components are of ground types. 

%% To expand the objects of aggregate types in a Firrtl
%% module, the process involves two phases. The first phase is to record
%% the types along with orientations and flipping information of the
%% objects, which helps determine the actual connection directions. For
%% example, if there is an input port of aggregate types with a field
%% marked as ``flip'', the ``flip'' mark will recursively apply to all
%% the atomic elements in that field. Since it is an input port, the
%% orientation is record as \source.  The second phase is to expand
%% the statements sequentially, and then produce a serial of new
%% statements.  For declaration statements, ground-typed elements are
%% generated corresponding to the original aggregate types and the
%% orientations. For connection statements, the connection directions are
%% determined by the orientations and flipping information from the first
%% phase. For subfields, there are two kinds of connections allowed for
%% HiFirrtl programs. One is connect subfields separately, in which case
%% the direction should already comply with the orientation of the
%% subfield elements. The other is to set connections between elements of
%% equivalent aggregate types. In such bundled connections, after
%% expansion process, individual subfield connections can be of different
%% directions.

In the sequel, we define how $\expandconnects(s)$ works for a statement $s$ in $M$. 
For each statement $s$, $\expandconnects(s)$ is a statement sequence. 
Then $\expandconnects(ss)$ for a statement sequence $ss$ can be easily defined as the sequential composition of $\expandconnects(s)$ for statements $s$ in $ss$. Finally, $\expandconnects(M)$ is defined as $\expandconnects(ss)$ where $ss$ is the sequence of all statements in $M$. 

%In the beginning of a Firrtl module, ports are declared first. 
For port declarations, 
when ports have flipped subfields, the expansion process will
change those fields of input ports into output ports, and vice
versa. %As a result, their orientations are changed between \source \, and \sink.
For an input port declaration statement $\expandconnects(\mbox{\prettyll{input} }x: t)$ is defined the following rule, 
 \begin{gather*}
   \label{expand-input}
   \infer[]
         {\seqq{\expandconnects(\mbox{\prettyll{input} }x: t) = pp}}
         {\begin{array}{c}
             {\listgexp(x, t)} = rl \hspace{15pt} \listgtyperef(t, \lfalse) = tl \\
             pp= \fmaptwo((fun\; (v, (gt, flipped)) \Rightarrow if\; flipped \;then\; \mbox{\prettyll{output}}\ v: gt  \;else\; \mbox{\prettyll{input} }v: gt), rl, tl)
           \end{array}
         }
 \end{gather*}
In the rule, the reference $x$ are expanded as a list $rl$ of ground-type components and the type $t$ are expanded as a list $tl$ of ground types.  Note that $rl$ and $tl$ have the same length, moreover, each element in $tl$ is of the form $(gt, flipped)$, where $gt$ is a ground type and $flipped$ is a Boolean value denoting the flipping status of $gt$. Then for each pair of elements in the same position $rl$ and $tl$, say $(v, (gt, flipped))$, if $flipped = \ltrue$, then a statement $\mbox{\prettyll{output}}\ v: gt$ is generated, otherwise, $\mbox{\prettyll{input} }v: gt)$ is generated. The function $\fmaptwo(f, l_1, l_2)$ used in the rule 
%Here, the {\tt map2}
is an operation that combines two lists $l_1, l_2$ of the same length into one list by traversing $l_1, l_2$ in the same pace and applying the binary function $f$ to each pair of elements that are currently visited. 
%The {\expandconnects} function takes a statement as input and
%produces a sequence of statements as output. In the case above, the
%input statement is the port declaration ``$\mbox{\prettyll{input }}x : t$'',
%and ``$pp$'' is the sequence of output statements. 
%The ``$pp$'' is produced by expanding the port reference and expanding its type.
%According to the flip information carried by the types, the input port may be flipped to a output port.

Expanding a register declaration is similar but slightly more involved. It is defined by the following rule, 
%
\begin{gather*}
  \label{expand-reg}
  \infer[]
        {\seqq{\expandconnects (\mbox{\prettyll{reg}}\; x:\,t\,clk\,\mbox{ \prettyll{with: reset=> }}(rsig, rval)) = ss}}
        {\begin{array}{c}
            \listgexp(x, t) = rl \hspace{15pt} \listgtyperef(t, \lfalse) = tl    \hspace{15pt}  \listgexp(rval, t) = rvall \\
            ss= \fmapthree((fun\; (v,(gt,\_),rv) \Rightarrow \mbox{\prettyll{reg}}\; v:\,gt\,clk \mbox{ \prettyll{with: reset=>}} (rsig, rv)), rl,tl,rvall)
          \end{array}
        }
\end{gather*}
Note that when expanding a register declaration, the reset values should also be expanded. Moreover, the $\fmapthree$ function is used to combine three lists $rl, tl, rvall$ into one list. 

Expanding a wire declaration is similar and simpler, thus omitted. 

%Here, {\expandconnects} function takes register statement
%``$\mbox{\prettyll{reg}}\; x:\,t\,clk\,\mbox{\prettyll{with: reset=>}} (rsig, rval)$''
%as input, which defines a register $x$. Similar to expanding a port,
%expanding a register involves expanding the reference to its subfield
%representations using {\expandref}, expanding the aggregate type
%using {\listgtyperef}. In addition, It includes the process of expanding
%the reset value with {\listgexp}.


%Elements such as wires and registers have duplex
%orientation. Regardless of the flip subfields in the original
%aggregate types, the resulting orientations will remain duplex. To
%expand registers in aggregate types, we also need to expand the reset
%values of the registers. Expanding wires is relatively simple and is
%omitted here.

Nodes are components of the {\source} orientation. 
They are only allowed to be defined with passive types, which means the type of the expression used in its definition should not have flipped fields, if there are fields in the type. 
If the type is not passive, then the expansion result is  an empty sequence. Expanding a node declaration is defined by the rule, 
\begin{gather*}
  \label{expand-node}
  \infer[]
        {\seqq{\expandconnects (\mbox{\prettyll{node }}\; x = e) = ss}}
        {\begin{array}{c}
            \typeofexp(e) = t  \hspace{15pt} \listgexp(x, t) = rl \hspace{15pt} \listgexp(e, t) = el \\
             ss= if\ \ispassive(t)\; then\;  \fmaptwo((fun\; (v,e) \Rightarrow \; \mbox{\prettyll{node}}\; v\;=\;e), rl,el)\ else\; []
          \end{array}
        }
\end{gather*}
where $\ispassive(t)$ is the function that checks whether $t$ contains flipped fields nor not. 
%In the definition above, function {\sf isPassive} recursively checks the the subfields of type $t$.
%If it has no flip field, it returns true; otherwise, it returns false.

To expand \prettyll{is invalid} statements, we
select only the {\sink} and {\duplex} elements to produce the sequence of
invalidations.
In HiFirrtl, it is allowed to have {\source} sub-elements in a bundle be
invalidated. The expansion result will discard them.
\begin{gather*}
  \label{expand-invalid}
  \infer[]
        {\seqq{\expandconnects(x\;\mbox{\prettyll{is invalid}}) = ss}}
        {\begin{array}{c}
             \typeofref(x) = (t, ori) \hspace{15pt} \listgexp(x, t) = rl \hspace{15pt}  \listgtyperef(t, \lfalse) = tl   \\
             ss = if(ori = \sink) \ then\  \fmaptwo(f_{\sink}, rl, tl)\ else\ if(ori = \source)\\
             \hspace{2.5cm} then\ \fmaptwo(f_{\source}, rl, tl)\ else\ \fmaptwo(f_{\duplex}, rl, tl)
          \end{array}
        } 
\end{gather*}
where the functions $f_{\sink},f_{\source}, f_{\duplex}$ are defined as follows, 
\[
\begin{array}{ l c l}
f_\sink & = & (fun\;(v, (gt, flipped)) \Rightarrow if\ \neg flipped \ then\ v\ \mbox{\prettyll{is invalid}}\ else\ [])
\\
f_\source & = & (fun\;(v, (gt, flipped)) \Rightarrow if\ flipped \ then\ v\ \mbox{\prettyll{is invalid}}\ else\ [])
\\
f_\duplex & = & (fun\;(v, (gt, flipped)) \Rightarrow v\ \mbox{\prettyll{is invalid}})
\end{array}
\]

Connections between bundled elements contain a set of connections for
individual subfields. According to the indicated ``flip'' for a subfield, the
connection direction may differ from the original.
\begin{gather*}
   \label{expand-cnct-ref}
   \infer[]
         {\seqq{\expandconnects (r\;\mbox{\prettyll{<=}}\;r') = ss}}
         {\begin{array}{c}
             \typeofref(r) = (t, ori) \hspace{15pt} \typeofref(r') = (t',ori') \\ 
             \listgexp(r, t) = rl \hspace{15pt} \listgexp(r', t') = rl'  \\
            \listgtyperef(t, \lfalse) = tl \hspace{15pt} \typecompatible(t, t', \lfalse) = b \\
             ss = if\;b\; then\; \fmapthree(fn, rl,rl',tl)\;else\,[]
           \end{array}
         }
\end{gather*}
where the function $fn$ is defined as follows,
\[
fn =  (fun\; (v_l,v_r,(\_,flipped)) \Rightarrow if\;flipped\; then\;v_r\, \mbox{\prettyll{ <= }}\, v_l\; else\;v_l\, \mbox{\prettyll{ <= }}\,v_r).
\]
%
In the case above, the connection is set between two references of
elements $r$ and $r'$. The function {\typecompatible} checks whether the
types can be connected. When the connection takes an
expression other than a reference on the right-hand side, flip
information is not taken into account. The expansion simply expands
the reference and the expression and produces a new sequence of
connections as follows.
%
\begin{gather*}
   \label{expand-cnct-expr}
   \infer[]
         {\seqq{\expandconnects (r\;\mbox{\prettyll{<=}}\;e) = ss}}
         {\begin{array}{c}
             \typeofref(r) = (t,\_) \hspace{15pt} \listgexp(r, t) = rl \hspace{15pt} \listgexp(e, t) = el \\
            \typecompatible(t,t', \lfalse) = b \hspace{15pt} ss = if\;b\; then\; \fmaptwo((fun\; (v,e) \Rightarrow v\, \mbox{\prettyll{ <= }}\, e), rl,el)\; else\;[]
           \end{array}
         }
\end{gather*}

%To summarize, a given Firrtl module can be expanded by applying the
%{\expandconnects} function to all its statements
%sequentially. Then the module in return contains only the elements and
%their connections in ground types.

